% Part: turing-machines
% Chapter: machines-computations
% Section: introduction

\documentclass[../../../include/open-logic-section]{subfiles}

\begin{document}

\olfileid{tur}{mac}{int}
\olsection{పరిచయం}

$\Nat$ నుండి $\Nat$కు వెళ్లే ఒక ప్రమేయం, ఉదాహరణకు,
\emph{గణనీయం} కావడం అంటే ఏమిటి? దీనికి మొదట వచ్చిన సమాధానాల్లో ఒకటి,
అలాగే అత్యంత ప్రసిద్ధమైనది: ఒక ప్రమేయాన్ని ట్యూరింగ్ యంత్రంతో గణించగలిగితే
అది గణనీయం. ఈ భావనను అలన్ ట్యూరింగ్ 1936లో ప్రతిపాదించాడు. ట్యూరింగ్
యంత్రాలు ఒక \emph{గణనా నమూనా}కు ఉదాహరణ---``గణనా ప్రక్రియ'' అనే భావనను
నిర్వచించడానికి అవి గణితపరంగా కచ్చితమైన పద్ధతి. దాని కచ్చితమైన అర్థం ఏమిటనే
విషయంపై భిన్నాభిప్రాయాలు ఉన్నా, గణనా ప్రక్రియలను నిర్దేశించడానికి ట్యూరింగ్
యంత్రాలు ఒక పద్ధతి అని విస్తృతంగా అంగీకరిస్తారు. ``ట్యూరింగ్ యంత్రం'' అనే పదం
కదిలే భాగాలున్న భౌతిక యంత్రపు బొమ్మను మనసులో రేకెత్తించినప్పటికీ, కచ్చితంగా
చెప్పాలంటే ట్యూరింగ్ యంత్రం పూర్తిగా ఒక గణిత నిర్మితి; అందువల్ల అది గణనా
ప్రక్రియ అనే భావనను ఆదర్శీకరిస్తుంది. ఉదాహరణకు, ట్యూరింగ్ యంత్రానికి పట్టే
కాలంపైనా జ్ఞాపకస్థలంపైనా మనం ఎలాంటి పరిమితీ విధించం: ఒక గణనకు ఈ విశ్వంలో
ఉన్న పరమాణువుల కంటే ఎక్కువ నిల్వ స్థలం లేదా ఎక్కువ అంచెలు అవసరమైనా, ట్యూరింగ్
యంత్రాలు దానిని గణించగలవు.

\begin{explain}
ఒక ప్రత్యేక రకమైన ఊహాత్మక యంత్రాంగానికి ప్రోగ్రాముగా ట్యూరింగ్ యంత్రాన్ని
భావించడం బహుశా ఉత్తమం. ఈ యంత్రాంగంలో ఒక \emph{టేపు}, ఒక
\emph{చదువు--వ్రాత శీర్షం} ఉంటాయి. మన ట్యూరింగ్ యంత్రాల రూపంలో టేపు ఒక
దిశలో (కుడివైపు) అనంతంగా ఉంటుంది; అది \emph{గడులు}గా విభజించబడి ఉంటుంది;
ప్రతి గడిలో ఒక పరిమిత \emph{వర్ణమాల}కు చెందిన ఒక సంకేతం ఉండవచ్చు. అలాంటి
వర్ణమాలల్లో ఎన్ని వేర్వేరు సంకేతాలైనా ఉండవచ్చు; కానీ ప్రధానంగా మనకు ఈ మూడే
చాలు: $\TMendtape$, $\TMblank$, $\TMstroke$. యంత్రాంగాన్ని ప్రారంభించినప్పుడు,
అత్యంత ఎడమ గడిలో $\TMendtape$, పరిమిత సంఖ్యలో కొన్ని గడుల్లో
\emph{ఇన్‌పుట్} ఉండటం మినహా టేపు ఖాళీగా ఉంటుంది (అంటే ప్రతి గడిలో
$\TMblank$ సంకేతం ఉంటుంది). ఏ సమయంలోనైనా యంత్రాంగం పరిమిత సంఖ్యలోని
\emph{స్థితులు}లో ఒకదానిలో ఉంటుంది. ప్రారంభంలో శీర్షం అత్యంత ఎడమ గడిని
చదువుతూ, నిర్దేశిత \emph{ప్రారంభ స్థితి}లో ఉంటుంది. యంత్రాంగపు నడకలోని ప్రతి
అంచెలో, ప్రస్తుతం చదువుతున్న గడిలోని విషయం, యంత్రాంగం ఉన్న స్థితి, ట్యూరింగ్
యంత్ర ప్రోగ్రాము---ఈ మూడూ కలిసి తరువాత ఏమి జరగాలో నిర్ణయిస్తాయి. ట్యూరింగ్
యంత్ర ప్రోగ్రామును ఒక పాక్షిక ప్రమేయం ఇస్తుంది; అది స్థితి~$q$, సంకేతం~$\sigma$
లను ఇన్‌పుట్‌గా తీసుకుని $\tuple{q', \sigma', D}$ అనే త్రయాన్ని అవుట్‌పుట్‌గా
ఇస్తుంది. యంత్రాంగం స్థితి~$q$లో ఉండి సంకేతం~$\sigma$ను చదివినప్పుడల్లా,
ప్రస్తుత గడిలోని సంకేతాన్ని $\sigma'$తో మార్చుతుంది; $D$ వరుసగా $\TMleft$,
$\TMright$, లేదా $\TMstay$ అయినదాన్ని బట్టి శీర్షం ఎడమకు, కుడికి కదులుతుంది
లేదా ఉన్నచోటే ఉంటుంది; యంత్రాంగం స్థితి~$q'$లోకి వెళుతుంది.

ఉదాహరణకు, \olref{fig:tm}లోని పరిస్థితిని పరిశీలించండి.
\begin{figure}
  \olasset{\olpath/assets/diagrams/turing-machine.tikz}
  \caption{తన ప్రోగ్రామును అమలు చేస్తున్న ట్యూరింగ్ యంత్రం.}
  \ollabel{fig:tm}
\end{figure}
ట్యూరింగ్ యంత్రపు టేపులో కనిపిస్తున్న భాగంలో అత్యంత ఎడమ గడిలో
టేపు-అంత సంకేతం $\TMendtape$ ఉంది; దాని తరువాత మూడు $1$లు, ఒక $0$, మరొక
నాలుగు $1$లు ఉన్నాయి. శీర్షం ఎడమ నుండి మూడవ గడిలోని $1$ను చదువుతూ
స్థితి~$q_1$లో ఉంది---దీనిని ``యంత్రం స్థితి~$q_1$లో ఒక $1$ను చదువుతోంది''
అంటాం. ఇన్‌పుట్ $\tuple{q_1, 1}$కు ట్యూరింగ్ యంత్ర ప్రోగ్రాము
$\tuple{q_2, 0, \TMstay}$ అనే త్రయాన్ని తిరిగి ఇస్తే, యంత్రం ఇప్పుడు మూడవ
గడిలోని $1$ను $0$తో మార్చి, చదువు--వ్రాత శీర్షాన్ని ఉన్నచోటే ఉంచి,
స్థితి~$q_2$కు మారుతుంది. తరువాత ఇన్‌పుట్ $\tuple{q_2, 0}$కు ప్రోగ్రాము
$\tuple{q_3, 0, \TMright}$ను తిరిగి ఇస్తే, యంత్రం ఆ $0$పైన మరొక $0$ను
వ్రాస్తుంది (అంటే చదువు--వ్రాత శీర్షం కింద ఉన్న టేపు విషయాన్ని వాస్తవానికి
మార్చదు), ఒక గడి కుడికి కదిలి స్థితి~$q_3$లోకి వెళుతుంది. ఇదే విధంగా
కొనసాగుతుంది.

ఏదైనా స్థితి $q_n$, సంకేతం $\sigma$ ఎదురై, $\tuple{q_n, \sigma}$కు ఎలాంటి
నిర్దేశమూ లేకపోతే, అంటే ఇన్‌పుట్ $\tuple{q_n,\sigma}$కు సంక్రమణ ప్రమేయం
నిర్వచితం కాకపోతే, యంత్రం \emph{ఆగుతుంది} అంటాం. మరో మాటలో చెప్పాలంటే,
యంత్రం అమలు చేయడానికి నిర్దేశమేదీ ఉండదు; ఆ బిందువులో దాని పని ఆగిపోతుంది.
ఆగడాన్ని కొన్నిసార్లు ప్రత్యేక ఆగే స్థితి~$h$తో సూచిస్తారు. దీనిని తరువాత
మరింత వివరంగా చూపిస్తాం.
\end{explain}

\begin{digress}
ట్యూరింగ్ రాసిన ``గణనీయ సంఖ్యల గురించి'' అనే వ్యాసపు అందం ఏమిటంటే, అతడు
ఒక ఆచారబద్ధ నిర్వచనాన్ని మాత్రమే కాక, ఆ నిర్వచనం గణనీయత అనే అంతర్బోధాత్మక
భావనను పట్టుకుంటుందని వాదనను కూడా అందించాడు. నిర్వచనం నుండి చూస్తే,
ట్యూరింగ్ యంత్రంతో గణించదగిన ఏ ప్రమేయమైనా అంతర్బోధాత్మక అర్థంలో గణనీయమే
అని స్పష్టంగా ఉండాలి. దీని విపర్యయం నిజమని---అంటే మనం సహజంగా గణనీయం అని
భావించే ఏ ప్రమేయాన్నైనా అలాంటి యంత్రంతో గణించవచ్చని---చూపడానికి ట్యూరింగ్
మూడు రకాల వాదనలను ఇచ్చాడు. అవి (ట్యూరింగ్ మాటల్లో):
\begin{enumerate}
\item అంతర్బోధను నేరుగా ఆశ్రయించడం.
\item రెండు నిర్వచనాల సమానత్వాన్ని నిరూపించడం (కొత్త నిర్వచనం
  అంతర్బోధాత్మకంగా మరింత ఆకర్షణీయమైతే).
\item గణనీయమైన సంఖ్యల విస్తృత వర్గాలకు ఉదాహరణలు ఇవ్వడం.
\end{enumerate}
మన లక్ష్యం గణనీయత భావనను ``సూత్రప్రాయంగా'' నిర్వచించడానికి ప్రయత్నించడం;
అంటే కాలం, స్థలం అనే ఆచరణాత్మక పరిమితులను పరిగణనలోకి తీసుకోకుండా
నిర్వచించడం. గణనీయతకు అత్యంత విస్తృత నిర్వచనం ఏర్పడిన తరువాత పరిమిత
వనరులతో చేసే గణనను పరిశీలించవచ్చు; ``గణనా సంక్లిష్టత''గా పిలిచే అధ్యయన
విషయానికి ఇదే కేంద్రం.
\end{digress}

\begin{history}
అలన్ ట్యూరింగ్ 1936లో ట్యూరింగ్ యంత్రాలను ఆవిష్కరించాడు. అప్పట్లో అతని
ఆసక్తి ప్రథమ-శ్రేణి తర్కం యొక్క నిర్ణేయతపై ఉన్నప్పటికీ, ఆ వ్యాసాన్ని కంప్యూటర్
రూపకల్పన పునాదులపై నిర్ణాయక వ్యాసంగా వర్ణించారు. ఆ వ్యాసంలో ట్యూరింగ్ గణనీయ
వాస్తవ సంఖ్యలపై---అంటే దశాంశ విస్తరణలను గణించగల వాస్తవ సంఖ్యలపై---దృష్టి
పెట్టాడు; అయితే తన భావనలను సహజ సంఖ్యలపై గణనీయ ప్రమేయాలకు, తదితరాలకు
అన్వయించడం కష్టం కాదని పేర్కొన్నాడు. మొదటి పనిచేసే సామాన్య-ప్రయోజన కంప్యూటర్‌ను
1941లో నిర్మించడానికి (జర్మనీకి చెందిన కొన్రాడ్ సూజె తన తల్లిదండ్రుల ఇంటి
గదిలో నిర్మించాడు) పూర్తి ఐదేళ్ల ముందే ఇది జరిగిందని గమనించండి. అలాగే
ట్యూరింగ్, బ్లెచ్లీ పార్క్‌లోని అతని సహచరులు సంకేతభేదక కొలాసస్‌ను నిర్మించడానికి
(1943) ఏడేళ్ల ముందు; అమెరికా ENIACకు (1945) తొమ్మిదేళ్ల ముందు; మొదటి బ్రిటిష్
సామాన్య-ప్రయోజన కంప్యూటర్---మాంచెస్టర్ స్మాల్-స్కేల్ ఎక్స్‌పెరిమెంటల్
మెషీన్---మాంచెస్టర్‌లో నిర్మించడానికి (1948) పన్నెండేళ్ల ముందు; అమెరికన్లు
BINACను మొదట పరీక్షించడానికి (1949) పదమూడేళ్ల ముందు ఇది జరిగింది. మాంచెస్టర్
SSEMకు మొదటి నిల్వ-ప్రోగ్రాము కంప్యూటర్ అనే ప్రత్యేకత ఉంది---అంతకుముందు
యంత్రాల్లో ప్రతి కొత్త పనికీ చేతితో తీగలను మళ్లీ అమర్చవలసి వచ్చేది.
\end{history}

\end{document}
